\documentclass[12pt]{article}
\usepackage[utf8]{inputenc}
\usepackage[T1]{fontenc}
\usepackage{hyperref}
\title{Dawawa Language Resources for K-12 Curriculum Development}
\author{Compiled by Claude Sonnet 5.5 (claude-sonnet-5-5) for Andr\'as Kornai}
\date{}
\begin{document}
\maketitle

\section{Introduction}
Dawawa (ISO 639-3: dww) is an Austronesian (Oceanic, Papuan Tip) language of Milne Bay Province, Papua New Guinea. Speaker estimates vary by source (about 2,000 in Wikipedia, 6,900 in Joshua Project citing Ethnologue). Very little was found online beyond a Bible translation and general descriptions.

\section{Primary Educational Resources}
No primers, textbooks or school programmes for Dawawa were found.

\section{Linguistic and Reference Materials}
Wikipedia overview: \url{https://en.wikipedia.org/wiki/Dawawa_language}. No grammar or dictionary was found in the searches.

\section{Religious and Community Resources}
\begin{itemize}
\item Dawawa New Testament (Vari Verenama), online at eBible: \url{https://ebible.org/dww/FRT01.htm}
\item Joshua Project people-group profile: \url{https://joshuaproject.net/people_groups/11515/PP}
\end{itemize}

\section{Digital and Audio Resources}
None identified beyond the online scripture text above.

\section{Curriculum Development Context}
The language is reported as stable and used by all in the community, but not known to be taught in schools. The scripture text is the only substantial written material found and could serve as a starting corpus, subject to community consultation.
\end{document}